\subsection{21.132: a graded radical quotient with finite centre}

The experiment ultimately classified this as a deduction from
prior radical theorems, rather than a new solution.
Start with a two-generated, positively graded, infinite-dimensional
nil \(\F_p\)-algebra \(A\), with finite-dimensional graded pieces,
from the graded Golod construction \cite{ershov-gs}.
Let \(L(A)\) be its Levitzki radical, the largest locally
nilpotent ideal. The centreless quotient \(B=A/L(A)\) is
the construction used by Sereda--Sozutov \cite{sereda-sozutov}.
The radical is proper, since otherwise finite generation
would make \(A\) nilpotent. It is homogeneous by the classical
graded-radical theorem explicitly recorded in
\cite{hong-nilradicals}. Hence \(B\) remains positively graded
with finite-dimensional pieces. It is infinite-dimensional:
otherwise positive grading makes it nilpotent, contradicting
\(L(B)=0\).

If \(z\) commutes with all of \(B\), then \(z\) is nilpotent
and the ideal \(B^1z\), using the unitization, is nilpotent.
It must be zero, so \(z=0\).
Let \(H=\langle1+b_1,1+b_2\rangle\le1+B\).
It is a \(p\)-group, since each of its elements is one plus
a nilpotent element. Its centre is trivial: a central
\(1+b\) commutes with the algebra generators, so \(b\) is
central in \(B\). It is infinite, since a finite \(H\) would
have finite-dimensional linear span containing the algebra
generators and hence all of \(B\).
The truncations \(B/B_{\ge n}\) give finite \(p\)-group
quotients separating every element, by finite homogeneous
support. Thus \(H\) is residually finite.
Adjoin a one-dimensional zero-multiplication algebra
\(\F_pz\). The subgroup generated by one plus the three
algebra generators is \(H\times C_p\), with centre \(C_p\).
It meets the Notebook's algebraic definition of a Golod
group, for every prime \(p\), and is at most three-generated.
The grading is what proves residual finiteness; it is not
inferred from an arbitrary quotient preserving that property.
The exact deduction under the question number was not located
in a prior paper, but its inputs and classification are explicit.
See \artifact{research/21.132-prior-consequence.md}.

\subsection{21.42: positive gradings and self-similar actions}

Every three-generated torsion-free nilpotent group of class
three is self-similar, giving a negative answer to the
existence question. Dekimpe--Igodt--Pouseele proved the
positive-grading assertion for the corresponding Lie
algebras \cite{dekimpe-igodt-pouseele}; the exact positive
form is restated in \cite{dekimpe-dere}.
Mathieu's criterion transfers it to the original lattice,
not merely an unspecified finite-index subgroup
\cite[Proposition~6]{mathieu}.
The rational Lie algebra of the Malcev completion is
generated by the logarithms of the three group generators,
so these results apply.

The mechanism is instructive. Positive weights give dilations
by \(q^j\) in weight \(j\). For a suitable integer \(q>1\)
the dilation \(\delta\) preserves the given lattice \(G\),
with finite-index image. Its inverse is a virtual
endomorphism \(\delta(G)\to G\).
If a subgroup \(K\subseteq\delta(G)\) is invariant under
this inverse, every \(\delta^{-n}(\log k)\) remains in the
discrete lattice and tends to zero. Thus \(k=1\).
The virtual endomorphism has trivial normal core and gives
a faithful recurrent action on the coset tree.
There is no prescribed degree or finite-state requirement
in the printed question.
The archive also retains an independent classification
argument giving positive weights for the three-generator
class-three Lie algebras, with weight choices
\((1,1,2)\), \((2,2,1)\), and \((3,4,2)\) in the
nonhomogeneous cases. That proof is a rediscovery of the
grading assertion; the lattice step remains an imported
theorem. See \artifact{research/21.42-known-consequence.md}.

\subsection{20.33: an effective universal group, with oracle bookkeeping}

For any oracle \(X\), there is a three-generated
\(X\)-computably-enumerably presented group \(G_X\) containing
every finitely generated group with an \(X\)-c.e. presentation.
This is a standard embedding consequence; Mikaelian's
explicit universal words improve three generators to two
\cite{mikaelian}.

Dovetail all oracle programs enumerating relators on each
finite alphabet, ignoring ill-formed outputs. Take the
free product \(U_X\) of all their presented groups.
It has an \(X\)-c.e. presentation on a countable computable
alphabet \(g_1,g_2,\ldots\), and each factor embeds.
For any such countably presented \(U\), in \(U*F(a,b)\)
the elements \(q_i=a^iba^{-i}\) freely generate a subgroup,
as part of a Schreier basis. The elements \(g_iq_i\) also
freely generate: the retraction killing \(U\) detects
any freely reduced relation among them.
The HNN extension
\[
 \langle U,a,b,t\mid t^{-1}q_it=g_iq_i\ (i\ge1)\rangle
\]
therefore embeds \(U\), and its relations express
\(g_i=t^{-1}q_itq_i^{-1}\).
Eliminating the \(g_i\) gives a three-generator presentation
by effective substitution in the old relators.
No word-problem decision is required, and the same oracle
suffices.

Conversely, a finitely generated subgroup of an \(X\)-c.e.
presented group has such a presentation: enumerate products
of conjugates of enumerated relators and compare their
freely reduced words with substitutions in the subgroup
generators. Dovetailing enumerates exactly the subgroup
kernel. Finite free products and adding finitely many
relators preserve \(X\)-c.e. presentability as well.
Thus both directions of the Notebook formulation follow;
in the forward direction take one copy of \(G_X\) and
the empty extra relator set. Decidable word problem and
finite presentation are not additional requirements.
See \artifact{research/20.33-known-consequence.md}.

\subsection{20.8 and 20.11: free-group deductions with the exact quantifiers}

For \kn{20.8}, the Lei--Zhang construction supplies a
counterexample \cite{lei-zhang}. In \(F(a,b)\) put
\[
 u=(a^{-1}b)^2,\quad v=(a^{-2}b^2)^2,\quad
 H=\langle a,u,v\rangle,\quad
 K=\langle u,aua^{-1},v,a^2va^{-2}\rangle.
\]
The ranks are \(3,4\), respectively. The homomorphism
on the free basis of \(H\) sending \(a\) to \(b\) and
fixing \(u,v\) differs from inclusion but fixes \(K\):
commutativity of \(c_i=a^{-i}b^i\) with \(c_i^2\) gives
\(a^ic_i^2a^{-i}=b^ic_i^2b^{-i}\).
To check the additional intermediate-subgroup condition,
identify \(H\) with \(F(t,x,y)\). The core graph of
\(\langle x,txt^{-1},y,t^2yt^{-2}\rangle\) has vertices
\(0,1,2\), a \(t\)-path between them, \(x\)-loops at
\(0,1\), and \(y\)-loops at \(0,2\).
Its rank is four. Of its five vertex partitions, no
identification and \(0=2\) give rank-four folded images;
the other three give the full three-loop rose.
The image in the core graph of any intermediate subgroup
is one of these. A rank-three image forces that subgroup
to be all of \(H\). Thus every proper intermediate
subgroup has rank at least four, including infinitely
generated ones. The ambient rank is two, which the
question permits. Full folded graphs are retained in
\artifact{research/20.8-known-consequence.md}.

For \kn{20.11}, Jaikin-Zapirain's prior rigidity theorem
\cite{jaikin} implies all four assertions, even when
the initial subgroup \(F\) and ambient free group \(H\)
are not finitely generated. Let \(\mathcal C\) be the
nonempty family of overgroups of \(F\) in \(H\) of least
finite rank \(r\).
For \(G_1,G_2\in\mathcal C\), \(G_1\) is compressed in
\(\langle G_1,G_2\rangle\), hence inert by Corollary~1.2.
Their intersection has rank at most \(r\), and thus exactly
\(r\) by minimality. Apply Corollary~1.5 to this now
finitely generated intersection to show that the join
also has rank \(r\).
A chain union has rank at most \(r\): otherwise \(r+1\)
members of a free basis generate a free factor lying
in one chain member, forcing its rank too large.
Zorn's lemma and pairwise joins give a greatest member.
For a least member, work inside one countable \(G_0\).
A countable subfamily has the same intersection as
all of \(\mathcal C\), by choosing a subgroup excluding
each element of \(G_0\) outside that intersection.
Finite intersections form a descending sequence of
rank \(r\). Any finite neighborhood in the intersection's
Schreier graph eventually appears unchanged in those
graphs. A finite subgraph of Betti number greater than
\(r\) is therefore impossible. The intersection has
rank \(r\) and is the least member. These limit steps
are needed to preserve the printed quantifiers.
See \artifact{research/20.11-known-consequence.md}.

\subsection{19.9, 16.14, and 19.48: further prior answers}

For \kn{19.9}, both answers were already explicit in
the pre-experiment GroupTheory50 preprint
\cite{proofforum-tensor}. The archive's alternative
deduction uses \(G=\operatorname{Sp}(10,\Z)\).
Its nonabelian tensor square is its universal central
extension, identified with the pullback of the real
topological universal cover in this stable rank.
Deligne's theorem makes that extension non-residually
finite \cite{deligne,benson-symplectic}.
It is finitely generated: lifts of generators generate
any perfect central extension of a finitely generated
group. Residual finiteness of finitely generated linear
groups then excludes linearity over every field.

For braid groups \(B_n\), \(n\ge4\), their abelianization
\(\Z\) splits the tensor square as
\(\Z\times(B_n\wedge B_n)\). The exterior square has
kernel \(H_2(B_n)=C_2\) over \(B_n'\)
\cite{bardakov-tensor,akita-braids}.
In the rational Clifford algebra with
\(e_i^2=1/2\) and \(e_ie_j=-e_je_i\), put \(u_i=e_i-e_{i+1}\).
Then \(u_i^2=1\), adjacent generators satisfy the braid
relation, and distant generators anticommute.
The induced map from the free presentation sends all
relators to central signs and detects the nontrivial
Hopf-kernel element \([x_1,x_3]\) as \(-1\).
Combining it with the commutator image embeds the exterior
square in \(B_n'\times A_n^\times\).
Left multiplication on the \(2^n\)-dimensional Clifford
algebra and the faithful Lawrence--Krammer representation
\cite{krammer} give a faithful tensor-square representation
of dimension \(1+\binom n2+2^n\).
This detects the central kernel explicitly; it does not
assume that every finite central extension of a linear
group is linear. See \artifact{research/19.9-proof.md}.

For \kn{16.14}, centrality of all involutions makes
\(E=\Omega_1(Z(G))=\Omega_1(G)\) a maximal elementary
abelian normal subgroup, of rank \(r=d(Z(G))\).
Sambale's earlier centralizer bound gives
\(|C_G(E):\Phi(C_G(E))|\le2^{2r}\) \cite{sambale-rank}.
Since \(C_G(E)=G\) and \(\Phi(G)=G^2\) for a finite
2-group, the requested inequality is
\(d(G/G^2)\le2d(Z(G))\). Products of quaternion groups
attain equality. This is a specialization of the prior
bound, not a new theorem; see
\artifact{research/16.14-prior-result.md}.

For \kn{19.48}, another pre-experiment GroupTheory50
preprint explicitly gives a nonclosed irreducible net
\cite{proofforum-net}. With \(R=\Z\), let \(\alpha\)
be a root of \(X^3-X-1\) and \(T=\Z[\alpha]\).
Its exceptional levels are
\[
 A=\Z+\Z\alpha+2\Z\alpha^2,\qquad
 B=\Z+2\Z\alpha+\Z(\alpha^2-1),
\]
with other levels \(2T\).
The identity \(\alpha^{-1}=\alpha^2-1\) allows root
transvections to produce the diagonal
\(\operatorname{diag}(\alpha,\alpha^{-1})\), whose
conjugation yields the forbidden upper parameter
\(\alpha^2\notin A\). The common ideal \(2T\) ensures
the net conditions and nonzero levels.
These are \(R\)-modules, as the question requires,
rather than subspaces over the fraction field.
The independently found characteristic-three example
remains in \artifact{research/19.48-proof.md}.
The earlier posted proofs are preprints; no outside
human acceptance is inferred from their availability.

\subsection{12.69: a literal counterexample retained as a formulation issue}

The printed condition permits a nonzero element of
\(\Z G\) sending each field element into a subfield.
For \(F=\Q(i)\), \(S=\Q\), and \(G=\{1,\sigma\}\)
with \(\sigma\) complex conjugation, the fixed element
\(1+\sigma\) acts multiplicatively as the norm
\(x\sigma(x)\in\Q\). Its nonnegative exponents also
give zero at \(x=0\). Additive interpretation gives
the trace and the same conclusion.
The extension is proper and separable.
The example \(F=\C,S=\R\) likewise contradicts the
printed claim that the uncountable analogue is affirmative.
This points to a missing qualification or other formulation
problem; changing from norm to trace does not repair it.
The experiment retained the literal observation, proposed
no invented correction, and counted no substantive resolution
of the intended question. See
\artifact{research/12.69-formulation-audit.md}.
